\section{Hasil}
\label{sec:results}

\subsection{Eksperimen A: Model yang melayani bukan model yang diminta}
Tabel~\ref{tab:served} merangkum uji substitusi. Model yang berhak dilayani sesuai dengan permintaan. Setiap nama yang tidak berhak---termasuk yang fiktif---diterima tanpa kesalahan dan dijawab oleh Claude Haiku 4.5. Tidak ada bidang dari objek respons yang berbeda antara kedua kasus; hanya event penggunaan (usage event) yang menamai model yang melayani. Agen yang menampung insiden tersebut telah memberi label beberapa bulan analisis dengan nama yang diminta; label-labelnya sekarang diturunkan dari event penggunaan (usage event), sebuah peringatan substitusi dicatat sekali per pasangan, dan endpoint kesehatan-nya melaporkan autentikasi secara eksplisit, karena login yang kadaluarsa sebelumnya hanya muncul sebagai hasil kosong.

\begin{table}[t]
  \caption{Eksperimen A: model yang diminta versus model yang melayani melalui SDK vendor pada akun paket dasar (berhak: Claude Haiku 4.5, GPT-5 mini, GPT-5.3 Codex, GPT-4.1). Tidak ada permintaan yang menghasilkan kesalahan.}
  \label{tab:served}
  \centering
  \footnotesize
  \begin{tabular}{@{}ll@{}}
    \toprule
    Diminta & Dilayani (event penggunaan (usage event)) \\
    \midrule
    Claude Haiku 4.5 & Claude Haiku 4.5 \\
    GPT-5 mini & GPT-5 mini \\
    Claude Opus 4.8 & Claude Haiku 4.5 \\
    nama fiktif & Claude Haiku 4.5 \\
    \bottomrule
  \end{tabular}
\end{table}

\subsection{Eksperimen B: Tiga Konfigurasi, Unit yang Sama}
Tabel~\ref{tab:results} dan Gambar~\ref{fig:configs} melaporkan perbandingan tersebut. Ketiga konfigurasi menghasilkan paper bahasa Indonesia yang dapat dikompilasi (sembilan halaman dibanding delapan dalam bahasa Inggris; nol referensi atau sitasi yang tidak terdefinisi; gerbang sitasi lolos dengan 32 kunci yang sama). Pada metrik kualitas otomatis, konfigurasi berada dalam satu galat baku satu sama lain: run tunggal Claude Haiku~4.5 memiliki mean cosine tertinggi (0.837 langsung, 0.961 terjemahan balik) dan kepatuhan glosarium (99.7\%), pipeline tiga-model mengikuti (0.826, 0.958, 99.0\%), dan model lokal berada di antaranya pada cosine (0.831, 0.956) tetapi terendah pada kepatuhan glosarium (97.9\%). Kedua konfigurasi model API mempertahankan semua 146 unit tanpa menggunakan model cadangan (fallback), sedangkan model lokal meninggalkan dua paragraf dalam bahasa Inggris setelah keluaran placeholder-nya gagal divalidasi dua kali. Waktu dinding berbeda dengan faktor tiga (22 vs 7 vs 6.4 menit) dan penggunaan API berbeda dengan faktor dua (32 vs 17 vs 0 panggilan). Teks akhir dari dua konfigurasi manapun memiliki kesamaan karakter 0.89--0.91, yaitu\ ketiga sistem menghasilkan bahasa Indonesia yang pada umumnya dapat dipertukarkan.

\begin{table*}[t]
  \caption{Eksperimen B pada 146 unit (4{,}673 kata). Nilai cosine adalah rata-rata dari teks akhir dengan galat baku; integritas = unit yang lolos validator tanpa kembali ke bahasa Inggris.}
  \label{tab:results}
  \centering
  \footnotesize
  \begin{tabular}{@{}lccc@{}}
    \toprule
    Metrik & \textsc{multi} & \textsc{single}-Claude Haiku & \textsc{single}-Gemma 3 \\
    \midrule
    Integritas (unit) & 146/146 & 146/146 & 144/146 \\
    Kepatuhan glosarium & 284/287 (99.0\%) & 286/287 (99.7\%) & 281/287 (97.9\%) \\
    Cosine EN$\leftrightarrow$ID & $0.826\pm0.011$ & $0.837\pm0.011$ & $0.831\pm0.011$ \\
    Cosine EN$\leftrightarrow$terjemahan balik & $0.958\pm0.006$ & $0.961\pm0.005$ & $0.956\pm0.007$ \\
    Unit yang tersisa dalam bahasa Inggris & 1 (klaim yang dikutip) & 0 & 2 \\
    Waktu dinding (s) & 1{,}334 & 425 & 384 \\
    Detik-model per unit (draf + sunting pasca) & 7.2 + 5.0 & 3.0 & 2.8 \\
    Panggilan API (model yang melayani) & 16 GPT-5 mini + 16 Claude Haiku 4.5 & 17 Claude Haiku 4.5 & 0 \\
    \bottomrule
  \end{tabular}
\end{table*}

\begin{figure*}[t]
  \centering
  \includegraphics[width=0.95\textwidth]{fig2_configs.pdf}
  \caption{Eksperimen B: kualitas teks akhir (cosine dengan galat baku), kepatuhan glosarium dan waktu dinding dari ketiga konfigurasi.}
  \label{fig:configs}
\end{figure*}

\paragraph{Apa yang dilakukan penyunting}
Pada \textsc{multi} penyunting mengubah 38 dari 146 unit (26\%); rasio suntingan rata-rata di semua unit adalah 0.017, dan pada unit yang diubah 0.067 (maksimum 0.625). Seleksi QA mempertahankan teks yang ditinjau untuk 139 unit dan draf untuk 7 (Gambar~\ref{fig:review}a); kedua kandidat berada dekat diagonal, yaitu\ penyuntingan pasca jarang mengubah skor otomatis. Catatan perubahan menjelaskan koreksi terminologi dan tata bahasa: \emph{inferred} yang di-render sebagai ``diturunkan'' (derived) dikoreksi menjadi ``disimpulkan''; \emph{show} yang di-render sebagai ``dibuktikan'' (proven) menjadi ``ditunjukkan''; sebuah calque bertanda hubung ``sembilan-aturan'' menjadi ``dengan sembilan aturan''; dan satu kalimat bahasa Inggris lengkap yang model pembuat draf tinggalkan tanpa terjemahan pada paragraf pembuka diterjemahkan. Tidak satupun dari hal ini terlihat pada cosine rata-rata.

\begin{figure*}[t]
  \centering
  \includegraphics[width=0.95\textwidth]{fig3_review.pdf}
  \caption{\textsc{multi} pipeline: (a) skor QA per unit dari draf dan dari teks pasca-sunting dengan kandidat yang dipilih; (b) distribusi rasio suntingan oleh penyunting atas 38 unit yang diubah.}
  \label{fig:review}
\end{figure*}

\subsection{Eksperimen C: Siapa yang Menemukan Kesalahan Mana}
Tabel~\ref{tab:errors} mengklasifikasikan kesalahan yang kami temukan menurut lapisan yang menangkapnya. Tiga kelas tertangkap oleh validator dan tidak oleh apa pun selain itu. (i) \emph{Placeholder dipindahkan}: pada satu run pilot, penyunting, saat ``merapikan'' format angka, memindahkan placeholder untuk \verb|\%| ke depan angka (``hanya \ph{3}58'' untuk ``only 58\%''); pemeriksaan multiset lulus, sehingga aturan menempel pada angka ditambahkan dan sekarang menolak keluaran semacam itu. (ii) \emph{Terjemahan balik yang menggemakan}: pada 37 panggilan untuk unit pendek (33 kandidat hasil penyuntingan pasca dan 4 kandidat draf) model QA lokal mengembalikan teks bahasa Indonesia tanpa perubahan alih-alih bahasa Inggris; karena terjemahan balik kemudian sama dengan kandidat, cosine terjemahan balik runtuh (19 unit turun lebih dari 0.10) dan selektor keliru memilih draf untuk 25 unit yang dua kandidatnya identik. Aturan gema mengubah ini menjadi panggilan yang gagal untuk dicoba ulang, setelah itu seleksi menjadi 139/7. (iii) \emph{Placeholder yang hilang}: model lokal sebagai penerjemah tunggal kehilangan placeholder dalam dua paragraf meskipun sudah dicoba ulang; tanpa model cadangan (fallback) ini tetap dalam bahasa Inggris dan dilaporkan, yang merupakan perilaku yang benar. Kesalahan yang tertangkap oleh model bersifat semantik: sebuah kalimat yang tidak diterjemahkan, istilah teknis yang salah (\emph{confounder} menjadi ``perusak'' (destroyer) pada run lokal dan ``perancu'' atau ``konfonder'' di tempat lain; ``stokhastisitas'' salah eja), sebuah sel tabel ``Union'' yang dibiarkan dalam bahasa Inggris dan sebuah judul yang dipangkas menjadi ``Kritikus lintas'' pada satu run Claude Haiku---ketiga terakhir ditemukan melalui pembacaan manusia, yang tetap menjadi lapisan terakhir.

\begin{table}[t]
  \caption{Eksperimen C: kelas kesalahan dan lapisan yang menangkapnya (P = model penyuntingan pasca, Q = model QA, V = validator deterministik, H = pembacaan manusia).}
  \label{tab:errors}
  \centering
  \small
  \begin{tabular}{@{}p{4.1cm}cp{2.6cm}@{}}
    \toprule
    Kesalahan & Tertangkap oleh & Kejadian \\
    \midrule
    Placeholder dipindahkan sebelum angka & V & 1 unit (pilot) \\
    Model QA menggemakan inputnya & V & 37 panggilan \\
    Placeholder yang hilang (model lokal) & V & 2 unit \\
    Kalimat tidak diterjemahkan & P & 1 unit \\
    Istilah teknis salah & P, H & 4 istilah \\
    Ejaan salah (``stokhastisitas'') & H & 1 unit \\
    Sel atau judul dibiarkan/terpotong & H & 2 unit (Claude Haiku) \\
    \bottomrule
  \end{tabular}
\end{table}

\subsection{Jawaban terhadap Pertanyaan Penelitian}
\textbf{RQ1.} Tidak: SDK vendor dapat menyajikan model yang berbeda dari yang diminta tanpa menghasilkan kesalahan apa pun, dan satu-satunya label yang benar adalah model yang disebutkan dalam event penggunaan; agen merekamnya untuk setiap panggilan. \textbf{RQ2.} Pipeline tiga-model tidak meningkatkan kualitas rata-rata---sebuah model kuat tunggal menyamai atau melampauinya dalam sepertiga waktu---dan integritas struktur merupakan sifat dari validator daripada dari pipeline (kedua konfigurasi API-model mencapai 146/146); yang ditambahkan oleh pipeline adalah koreksi terhadap kesalahan semantik yang tidak ditunjukkan oleh metrik mana pun. \textbf{RQ3.} Kesalahan struktural dan protokol (placeholder, gema) hanya ditangkap oleh validator; kesalahan semantik ditangkap oleh model kedua atau oleh peninjau manusia.
